\section{Chapter 1 supplementary exercises}
These exercises combine consistency, parametrization, independence,
and transformations. Keep existence and uniqueness separate: a missing
variable pivot creates freedom only in a \emph{consistent} system.
When a problem describes physical quantities such as floors, solve the
real linear equations first, then impose the required integer and
nonnegative restrictions.

\exercise{1.S.1}{Twenty-six statements and their logical pitfalls}
\begin{enumerate}[label=(\alph*)]
\item \textbf{False.} Reduced echelon form is unique; ordinary echelon
form is not. The one-by-one matrices $[1]$ and $[2]$ are row equivalent
and both in echelon form, but are different.
\item \textbf{False.} For $n=1$, the equation $0x=0$ has infinitely
many solutions, more than one. A square system can have zero, one,
or infinitely many solutions.
\item \textbf{True.} If $Ax=Ay=b$ and $x\ne y$, then
$A(x+t(y-x))=b$ for every real $t$. Different $t$ give distinct
solutions because $y-x\ne0$.
\item \textbf{False.} No free variables guarantees at most one solution,
not consistency. The equations $x=0$ and $x=1$ have no free variable
and no solution.
\item \textbf{True.} Elementary row operations are reversible, so each
operation preserves the solution set in both directions.
\item \textbf{True.} If $Ax=Ay=b$ with $x\ne y$, then
$A(x-y)=0$ with $x-y\ne0$. Thus the homogeneous system has
both zero and a distinct nonzero solution.
\item \textbf{False.} Every $A$ has a consistent system for $b=0$.
For the zero $2\times2$ matrix the columns span only $\{0\}$, not $\R^2$.
\item \textbf{False.} Every matrix has a reduced echelon form, including
inconsistent augmented matrices. The single equation $0x=1$ is already
a counterexample.
\item \textbf{True.} Each row-equivalence class has one reduced echelon
form. Reducing either member therefore gives the same matrix.
\item \textbf{False.} $Ax=0$ always has the trivial solution, including
when free variables exist. Take $A=[0]$.
\item \textbf{True.} Consistency for every $b\in\R^m$ means a pivot
in every one of the $m$ rows. There are exactly $m$ pivots and thus
$m$ pivot columns.
\item \textbf{False.} Row pivots imply existence for every target,
not uniqueness. The matrix $[1\ 0]$ has a pivot in every row, but
$x_1=b$, with $x_2$ arbitrary, gives infinitely many solutions.
\item \textbf{True.} In a square $n\times n$ matrix, $n$ pivots
occupy every column and every row. Reduced pivots are ones with zeros
elsewhere in their columns, so the reduced matrix is $I_n$.
\item \textbf{True.} Both matrices reduce to $I_3$. Reduce $A$ to
$I_3$, then reverse the operations reducing $B$ to $I_3$ to reach $B$.
\item \textbf{True.} The first system supplies a nonzero $z$ with
$Az=0$. If $Ay=c$, then $A(y+tz)=c$ for every real $t$,
giving infinitely many distinct solutions of the second system.
\item \textbf{True.} Row-equivalent matrices have the same number
of pivots. If $A$ spans $\R^m$, its $m$ pivots fill all rows;
the same is true of $B$, so its columns also span $\R^m$.
This does not assert equality of their actual column vectors.
\item \textbf{False.} The list $e_1,e_2,e_1+e_2$ in $\R^3$ has
no proportional pair but has the nontrivial relation
$e_1+e_2-(e_1+e_2)=0$.
\item \textbf{True.} Three vectors in $\R^2$ cannot be independent:
a $2\times3$ matrix has at most two pivots. Thus the assumed
independence excludes that ambient space.
\item \textbf{False.} Four columns supply at most four pivots, not
the five row pivots required to span $\R^5$.
\item \textbf{True.} $-u=(-1)u+0v$, a permitted linear combination.
\item \textbf{False.} Let $u=e_1$, $v=2e_1$, $w=e_2$ in $\R^2$.
All are nonzero, but the first two span only the horizontal axis,
which does not contain $w$.
\item \textbf{False.} Take $u=e_1$, $v=e_2$, and $w=e_2$.
Then $w=0u+v$, but $u$ is outside $\Span\{v,w\}$.
\item \textbf{False.} The hypothesis does not exclude $v_1=0$.
Take $v_1=0$, $v_2=e_1$, $v_3=e_2$ in $\R^5$. The second
is not a multiple of the first, and the third is outside their span,
but the zero vector makes the list dependent.
\item \textbf{True.} A linear transformation is a function with
additional operation-preserving properties.
\item \textbf{True.} A $6\times5$ matrix cannot have six pivots,
so its columns cannot span $\R^6$.
\item \textbf{False.} $m$ pivot columns mean onto $\R^m$.
For one-to-one we need $n$ pivots. Again $[1\ 0]$ has $m=1$
pivot but sends every $(0,t)^T$ to zero.
\end{enumerate}

\exercise{1.S.2}{Every possible solution set of one equation}
If $a\ne0$, division gives the unique solution $x=b/a$.
If $a=0$ but $b\ne0$, the equation says $0=b$ and has no solution.
If $a=b=0$, the equation is true for every real $x$, so its solution
set is all of $\R$. These three cases exhaust the values of $a,b$.

\exercise{1.S.3}{Three configurations of planes}
\begin{enumerate}[label=(\alph*)]
\item Use $x=0$, $y=0$, $x+y=0$. The first two force $x=y=0$,
while $z$ is arbitrary; the third adds no restriction. The intersection
is the $z$-axis, a single line.
\item Use $x=0$, $y=0$, $z=0$. Their only common point is the origin.
\item Use $x=0$, $x=1$, $y=0$. No point can satisfy the first two
simultaneously, so there is no common intersection.
\end{enumerate}
Every listed equation has at least one nonzero coefficient and therefore
really does describe a plane, rather than the whole space or an empty set.

\exercise{1.S.4}{Why the square size matters}
Three pivot columns in a $3\times3$ coefficient matrix also mean
three pivot rows. Thus appending any right-hand side gives a consistent
system, because no zero-coefficient row remains to create a contradiction.
There are no free variables, so the solution is unique. Both properties
are needed: column pivots alone in a taller matrix would not ensure existence.

\exercise{1.S.5}{Classify two parameter-dependent systems}
\step{(a) Eliminate the first variable.}
Subtract four times $x_1+3x_2=k$ from the second equation. The
result is $(h-12)x_2=8-4k$. If $h\ne12$, this determines $x_2$
uniquely, then the first equation determines $x_1$. If $h=12$,
the result is $0=8-4k$, consistent exactly when $k=2$.
\answer{(i) No solutions: $h=12,k\ne2$. (ii) One solution:
$h\ne12$ with any $k$. (iii) Infinitely many: $h=12,k=2$.}
\step{(b) Add three times the first equation to the second.}
The eliminated equation is $(k+3h)x_2=1$.
If $k+3h\ne0$, then $x_2=1/(k+3h)$ and
$x_1=(hx_2-1)/2$ are uniquely determined. If $k+3h=0$,
we obtain the contradiction $0=1$.
\answer{(i) No solutions: $k=-3h$. (ii) One solution:
$k\ne-3h$. (iii) No parameter values give infinitely many solutions.}

\exercise{1.S.6}{Restate one consistency question three ways}
Let $a_1=(4,8)^T$, $a_2=(-2,-3)^T$, $a_3=(7,10)^T$,
and $b=(-5,-3)^T$.
\step{(a) Use a linear combination.}
Ask whether $b=x_1a_1+x_2a_2+x_3a_3$ for some real weights.
Eliminating twice the first scalar equation from the second gives
$x_2-4x_3=7$. Let $x_3=t$; then $x_2=7+4t$ and
$4x_1=9+t$. Thus the full solution family is
\[
 x=(9/4+t/4,\ 7+4t,\ t)^T,\qquad t\in\R.
\]
In particular $(9/4,7,0)^T$ works, proving consistency.
\step{(b) Use matrix columns.}
For $A=\begin{bmatrix}4&-2&7\\8&-3&10\end{bmatrix}$, ask
whether $b$ lies in the span of the columns of $A$. Part (a) answers yes.
\step{(c) Use a transformation.}
Define $T:\R^3\to\R^2$ by $T(x)=Ax$. Ask whether $b$ is
in the range of $T$, or equivalently whether any input has image $b$.
The displayed family supplies all such inputs.

\exercise{1.S.7}{Restate consistency for every right-hand side}
Set $v_1=(2,-5,7)^T$, $v_2=(-4,1,-5)^T$, and
$v_3=(-2,1,-3)^T$, and put $A=[v_1\ v_2\ v_3]$.
\step{(a) Ask whether the span is all of $\R^3$.}
Row three of $A$ is row one minus row two. Hence every vector
$b$ in the column span satisfies $b_3=b_1-b_2$. The first two
rows have two independent leading columns: elimination of $x_1$
from them leaves a nonzero coefficient $-9$ on $x_2$ after
multiplying the first row by $5/2$ and adding it to row two.
Thus the first two equations can be solved for any $b_1,b_2$,
but the third requires exactly $b_3=b_1-b_2$.
Consequently $\Span\{v_1,v_2,v_3\}$ is this plane, not $\R^3$.
\step{(b) State the column version.}
The question asks whether the columns of $A$ span $\R^3$.
They do not; for example $(0,0,1)^T$ is outside their span.
\step{(c) State the transformation version.}
For $T:\R^3\to\R^3$, $T(x)=Ax$, the question is whether
$T$ is onto. It is not, because its range is the proper plane above.
\answer{The system is consistent precisely for $b_3=b_1-b_2$,
so it is not consistent for all triples of right-hand sides.}

\exercise{1.S.8}{Echelon patterns for spanning columns}
Use $p,q,r$ for nonzero pivots and $*$ for arbitrary entries.
\step{(a) Two rows and three columns.}
Spanning $\R^2$ requires two row pivots, which can occur in columns
$(1,2)$, $(1,3)$, or $(2,3)$:
\[
 \begin{bmatrix}p&*&*\\0&q&*\end{bmatrix},\quad
 \begin{bmatrix}p&*&*\\0&0&q\end{bmatrix},\quad
 \begin{bmatrix}0&p&*\\0&0&q\end{bmatrix}.
\]
\step{(b) Three rows and three columns.}
All three columns must also be pivot columns, leaving the sole pattern
$\begin{bmatrix}p&*&*\\0&q&*\\0&0&r\end{bmatrix}$.

\exercise{1.S.9}{Decompose a vector along two different lines}
A vector on $y=2x$ is $(s,2s)$; one on $y=x/2$ is $(t,t/2)$.
Their sum equals $(5,6)$ when
$s+t=5$ and $2s+t/2=6$. Substituting $t=5-s$ into the latter
and multiplying by two gives $4s+5-s=12$, so $s=7/3$ and
$t=8/3$. Therefore
\[
 (5,6)=(7/3,14/3)+(8/3,4/3).
\]
Each summand satisfies its respective line equation, and their
coordinate sums check the requested result.

\exercise{1.S.10}{Infer existence and uniqueness from the diagram}
The two source arrows $a_1,a_2$ are nonzero and not parallel,
so they are independent. Their two column directions span $\R^2$.
Thus any $b$ shown in that plane is attained, and its two coefficients
are unique. Negative coefficients are allowed; being outside the wedge
of positive combinations does not prevent a solution of $Ax=b$.
\answer{Yes, there is exactly one solution.}

\exercise{1.S.11}{A nonechelon matrix with a line of zero-images}
Choose
\[
 A=\begin{bmatrix}1&1&0\\1&0&0\end{bmatrix}.
\]
Both rows lead in column one, so it is not in echelon form.
The second equation in $Ax=0$ gives $x_1=0$, then the first
gives $x_2=0$, while $x_3$ is free. Thus the solution set is
$\{(0,0,t)^T:t\in\R\}$, the required line in $\R^3$.

\exercise{1.S.12}{A nonechelon matrix with a plane of zero-images}
Choose $A=\begin{bmatrix}1&2&3\\2&4&6\end{bmatrix}$.
The second row again leads under the first, so the matrix is not
echelon. Its second equation is twice the first, leaving only
$x_1+2x_2+3x_3=0$. Its solutions are
\[
 x=s(-2,1,0)^T+t(-3,0,1)^T,\qquad s,t\in\R.
\]
The two generators are independent, so this is a plane, not merely a line.

\exercise{1.S.13}{A prescribed null vector determines the missing entries}
The pivot information forces reduced echelon form
$\begin{bmatrix}1&0&u\\0&1&v\\0&0&0\end{bmatrix}$.
Applying it to $(3,-2,1)^T$ gives $(3+u,-2+v,0)^T$.
For this to be zero, $u=-3$ and $v=2$. Hence the required form is
\[
 \begin{bmatrix}1&0&-3\\0&1&2\\0&0&0\end{bmatrix}.
\]
Row operations preserve the null space, so the prescribed vector must
satisfy these reduced equations even if the original $A$ was different.

\exercise{1.S.14}{Parameter values for two-vector independence}
Put the two vectors into columns and replace $R_2$ by $R_2-aR_1$:
\[
 \begin{bmatrix}1&a\\a&a+2\end{bmatrix}
 \sim\begin{bmatrix}1&a\\0&a+2-a^2\end{bmatrix}.
\]
There are two pivots precisely when $a+2-a^2\ne0$.
Since $a+2-a^2=-(a-2)(a+1)$, independence holds for
$\boxed{a\ne2,-1}$. At either excluded value there is one pivot,
so a free coefficient supplies dependence.

\exercise{1.S.15}{Parameters in triangular-looking vector lists}
\step{(a) Three vectors in $\R^3$.}
The column matrix is $\begin{bmatrix}a&b&d\\0&c&e\\0&0&f\end{bmatrix}$.
Independence requires and is guaranteed by $a,c,f\ne0$; the other
entries $b,d,e$ may be any real numbers. For sufficiency, read a
homogeneous relation from bottom to top: $fz=0$ gives $z=0$,
then $cy=0$ gives $y=0$, then $ax=0$ gives $x=0$.
For necessity, if $a=0$ the first vector is zero; if $c=0$, the
first two vectors are dependent; if $f=0$, all three vectors lie in
the two-dimensional coordinate plane with third coordinate zero.
\step{(b) Three vectors in $\R^4$.}
There are \emph{no restrictions} on $a,b,c,d,e,f$.
A relation among the columns, read from the last coordinate upward,
first forces the third coefficient to zero because that coordinate
is $1$ only in column three. The third coordinate then forces the
second coefficient to zero, and the second coordinate forces the
first to zero. Hence the list is independent for every choice of the six numbers.

\exercise{1.S.16}{Use the preceding-column characterization}
The supplied matrix has nonzero diagonal entries $1,5,8,10$
and zeros above the diagonal. Column one is nonzero. For column
$j>1$, suppose it were a combination of the preceding columns.
Its first $j-1$ coordinates are zero. Reading those coordinates
from top to bottom forces every coefficient of the preceding columns
to be zero, because each relevant diagonal entry is nonzero.
The proposed combination would then be zero, contradicting the nonzero
$j$th coordinate of column $j$. Thus no later column is a combination
of preceding ones. The characterization in Theorem 7 of Section 1.7
now proves that all four columns are independent.

\exercise{1.S.17}{Adjoin a vector outside the old span}
Suppose $c_1v_1+c_2v_2+c_3v_3+c_4v_4=0$.
If $c_4\ne0$, solving for $v_4$ would put it in
$\Span\{v_1,v_2,v_3\}$, contrary to the hypothesis. So $c_4=0$.
Independence of the first three then forces $c_1=c_2=c_3=0$.
All coefficients vanish, proving independence of the four-vector list.

\exercise{1.S.18}{Add one basis vector to another}
A relation $av_1+b(v_1+v_2)=0$ can be collected as
$(a+b)v_1+bv_2=0$. Independence of $v_1,v_2$ forces
$b=0$ and $a+b=0$, hence $a=0$. Thus the new pair
$v_1,v_1+v_2$ is also independent. The argument proves the result
for all coefficients, not just a selected numerical example.

\exercise{1.S.19}{Three collinear points give dependent position vectors}
Because the points are distinct, $v_2-v_1\ne0$ gives a direction
for their common line. Some real $t$ satisfies
$v_3=v_1+t(v_2-v_1)=(1-t)v_1+tv_2$. Therefore
\[
 (1-t)v_1+tv_2-v_3=0.
\]
The coefficient of $v_3$ is $-1$, so this relation is nontrivial.
It proves dependence even if the line does not pass through the origin:
the line's position vectors still lie in the span of two of its points.

\exercise{1.S.20}{Reverse an input using homogeneity}
Since $-u=(-1)u$, linearity gives
$T(-u)=(-1)T(u)=-v$. This uses the scaling rule with the scalar
$-1$ and does not require $T$ to be invertible.

\exercise{1.S.21}{A coordinate-plane reflection}
The images of $e_1,e_2,e_3$ are $e_1,-e_2,e_3$.
Put them into columns to get
\[
 A=\begin{bmatrix}1&0&0\\0&-1&0\\0&0&1\end{bmatrix}.
\]
It leaves the first and third coordinates fixed and reverses the second,
exactly the specified reflection through $x_2=0$.

\exercise{1.S.22}{Surjectivity forces injectivity for this square size}
Onto $\R^3$ means a pivot in every row of the $3\times3$ matrix.
These three pivots also occupy every column, so there is no free
variable in $Ax=0$. Its only solution is zero, giving one-to-one.
The equality of the input and output dimensions is crucial here.

\exercise{1.S.23}{A Givens rotation in two coordinates}
Multiplying the displayed matrix by $(4,3)^T$ gives
$(4a-3b,4b+3a)^T$. The target requires
$4a-3b=5$ and $3a+4b=0$. From the second equation,
$b=-3a/4$; substituting in the first gives $25a/4=5$.
Thus $\boxed{a=4/5,b=-3/5}$.
The condition $a^2+b^2=16/25+9/25=1$ is satisfied, and the
image is $(5,0)^T$, so both direction and length check.

\exercise{1.S.24}{Zero one coordinate without changing the middle one}
The first and third coordinates give
$2a-4b=2\sqrt5$ and $2b+4a=0$. The latter says $b=-2a$,
so $10a=2\sqrt5$. Therefore
\[
 a=1/\sqrt5,\qquad b=-2/\sqrt5.
\]
The middle coordinate remains $3$ automatically, and
$a^2+b^2=1/5+4/5=1$. Substituting in the first and third
coordinates gives $2\sqrt5$ and $0$, as required.

\exercise{1.S.25}{A building design needs whole, nonnegative floor counts}
\step{(a) Interpret a scalar multiple of the recipe column.}
The vector $x_1(3,7,8)^T$ gives the numbers of three-, two-, and
one-bedroom apartments on the $x_1$ floors using plan A.
\step{(b) Add the contributions from all three plans.}
The total apartment-count vector is
\[
 x_1\begin{bmatrix}3\\7\\8\end{bmatrix}
 +x_2\begin{bmatrix}4\\4\\8\end{bmatrix}
 +x_3\begin{bmatrix}5\\3\\9\end{bmatrix}.
\]
\step{(c) Solve the target equations, then impose the physical restrictions.}
Appending $(66,74,136)^T$ and reducing gives
\[
 \left[\begin{array}{ccc|c}3&4&5&66\\7&4&3&74\\8&8&9&136\end{array}\right]
 \sim
 \left[\begin{array}{ccc|c}1&0&-1/2&2\\0&1&13/8&15\\0&0&0&0\end{array}\right].
\]
Thus the real solutions are $x_1=2+t/2$, $x_2=15-13t/8$,
$x_3=t$. Floors must be nonnegative integers. Integer $t$ makes
$x_2$ an integer only when $t$ is a multiple of $8$, since $13$
and $8$ have no common factor. Nonnegativity requires
$0\le t\le120/13<10$. Hence the only allowed values are $t=0,8$,
giving $(x_1,x_2,x_3)=(2,15,0)$ or $(6,2,8)$.
Substitution yields the requested apartment totals in both cases.
\answer{Yes. There are exactly two designs with whole nonnegative
floor counts, although the unrestricted real system has infinitely many solutions.}
